\begin{abstract}
Synchronous updates to independent files still contend for a few shared append states inside the filesystem.
On Zoned Namespace (ZNS) solid-state drives (SSDs), distributing appends must respect active and open zone limits while preserving synchronization semantics for all pending mappings.
These mappings include obligations from earlier asynchronous writeback, even after the corresponding data pages become clean.

Z-ParaPack organizes file-update commit closures together with append permissions.
An update pack (PACK) gathers current writeback and pending mappings; a terminal record in the last selected node defines durable completion.
Reusable append lanes separate allocation state across concurrent updates and reserve a lane for maintenance writes.
Recovery validates and replays complete PACKs in commit-number order.
On Linux 7.2.9, Z-ParaPack exceeds the fastest baseline under highly concurrent synchronous random writes and increases completion rates for mixed file operations, database transactions, and persistent messaging.
Both SQLite configurations improve transaction completion rates and achieve the lowest p99.9 among the evaluated systems.
Path analysis shows shorter shared-allocation intervals together with shorter complete synchronization calls at high concurrency.
\end{abstract}
\begin{IEEEkeywords}
ZNS SSD, file systems, file synchronization, parallel append, crash recovery.
\end{IEEEkeywords}
